%------------------------
% CareerOS master resume template
%
% This is Samarth's real resume with four substitution markers added. Everything outside the markers is copied
% through verbatim on every tailoring run, so the layout, packages and hand-written sections (Publications,
% Skills table, Leadership, Education) stay exactly as authored.
%
% Four markers are filled from data/master/achievements.yaml by a validated selection. Each is the marker name
% below wrapped in double percent signs, and each must appear EXACTLY ONCE in this file -- which is why they
% are named indirectly here rather than written out. (Writing one literally in a comment makes substitution
% happen here too, and since experience blocks span several lines, everything after the first line escapes the
% comment and duplicates a whole section into the document.)
%   CAREEROS:SUMMARY      one tailored paragraph positioning you for the specific role
%   CAREEROS:SKILLS       the subset of your skills most relevant to this job, ordered
%   CAREEROS:EXPERIENCE   selected work-experience entries and bullets
%   CAREEROS:PROJECTS     selected project entries and bullets
%
% Deliberately unchanged from the original except:
%   - the four markers replace the previously static Work Experience / Projects lists
%   - a Summary section was added (it was commented out in the original) since a per-job summary is one of the
%     highest-value things tailoring can do
%   - a "Most relevant for this role" line was added above the Skills table rather than replacing it: the
%     hand-written categorised table is better than a flat generated list, so it survives untouched
%   - the biblatex package and \addbibresource were removed. Publications are listed manually below, so
%     biblatex produced no output while making the document fail to build without own-bib.bib present.
%
% The original, untouched file is kept at the repository root as resume-current.tex.
%------------------------
\documentclass{article}
\usepackage{url}
\usepackage{parskip}

%other packages for formatting
\RequirePackage{color}
\usepackage{soul}

\RequirePackage{graphicx}
\usepackage[usenames,dvipsnames]{xcolor}
\usepackage[scale=0.9]{geometry}

%tabularx environment
\usepackage{tabularx}

%for lists within the experience section
\usepackage{enumitem}

% centered version of 'X' col. type
\newcolumntype{C}{>{\centering\arraybackslash}X}

%to prevent spillover of tabular into next pages
\usepackage{supertabular}
\usepackage{tabularx}
\newlength{\fullcollw}
\setlength{\fullcollw}{0.47\textwidth}

%custom \section
\usepackage{titlesec}
\usepackage{multicol}
\usepackage{multirow}

%http://stefano.italians.nl/archives/26
\titleformat{\section}{\large\scshape\raggedright}{}{0em}{}[\titlerule]
\titlespacing{\section}{0pt}{4pt}{4pt}

%Setup hyperref package, and colours for links
\usepackage[unicode, draft=false]{hyperref}
\definecolor{linkcolour}{rgb}{0,0.2,0.6}
\hypersetup{colorlinks,breaklinks,urlcolor=linkcolour,linkcolor=linkcolour}

%for social icons
\usepackage{fontawesome5}

\begin{document}

\pagestyle{empty}
\begin{center}
    \textbf{\Huge{Samarth Shekhar Thosar}} \\[8pt]
    \faMobile* \ \href{tel:+917887911070}{+91 7887911070} \textbar\
    \faEnvelope \ \href{mailto:samarth.thosar194@gmail.com}{samarth.thosar194@gmail.com} \textbar\
    \faLinkedin \ \href{https://linkedin.com/in/samarth-thosar}{samarth-thosar} \textbar\
    \faGithub \ \href{https://github.com/samarth-thosar}{samarth-thosar} \textbar\
    \faGraduationCap \ \href{https://scholar.google.com/citations?user=GROqtyIAAAAJ&hl}{Google Scholar}
\end{center}

\section{Profile Summary}
%%CAREEROS:SUMMARY%%

\vspace{10pt}
\section{Work Experience}
\vspace{4pt}
%%CAREEROS:EXPERIENCE%%

\section{Publications}
\begin{enumerate}

    \item \textbf{S. Thosar}, R. Bharadwaj, S. Ratnaparkhi, R. Rajpurohit, K. Rahate, and R. Pandita,
    “\textit{Deepfake Detection for Preventing Audio and Video Frauds using Advanced Deep Learning Techniques},”
    in \textit{Proc. of 2023 International Conference on Integrated Intelligence and Communication Systems (ICIICIS)},
    IEEE, Nov. 2023, pp. 1–7.
    DOI: \url{https://doi.org/10.1109/ICAICCIT60255.2023.10465762}
    \textit{(Indexed in Scopus)}

    \item \textbf{S. Thosar}, K. Ingale, S. Vaze, V. Uttarwar, S. Urane, and T. Mahajan,
    “\textit{Battlefield Health Surveillance with Soldier Tracking: A Privacy-Preserving Encrypted Approach},”
    in \textit{Proc. of 2023 International Conference on Advances in Computation, Communication and Information Technology (ICAICCIT)},
    IEEE, Nov. 2023, pp. 896–903.
    DOI: \url{https://doi.org/10.1109/ICIICS59993.2023.10421486}
    \textit{(Indexed in Scopus)}

\end{enumerate}


\section{Skills Summary}
\small
\textbf{Most relevant for this role:} %%CAREEROS:SKILLS%%
\vspace{4pt}

\renewcommand{\arraystretch}{0.8}
\begin{tabular}{p{4.8cm}p{13.5cm}}
  \textbf{Programming Languages:} & Python, C++, Java, C, TypeScript, SQL, Dart, R, JavaScript, HTML/CSS, LaTeX \\
  \textbf{Frameworks:} & React Native, ReactJS, Flutter, Node.js, Express.js, Redux, Spring Boot, Django \\
  \textbf{Libraries:} & TensorFlow, PyTorch, OpenCV, Scikit-learn, NumPy, Pandas, Matplotlib, NLTK \\
  \textbf{DevOps \& Deployment:} & Docker, Kubernetes, Jenkins, MLFlow, Git, GitHub, CI/CD Pipelines, Firebase, AWS, GCP \\
  \textbf{Databases:} & PostgreSQL, MySQL, MongoDB, Realm, Firebase Firestore \\
  \textbf{Tools \& Platforms:} & Figma, JIRA, Anthropic Claude API, MCP, Claude Code, Claude Cowork, Agent Skills \\
  \textbf{Core Competencies:} & Problem Solving, Leadership, Team Management, Research Communication, Public Speaking, Collaboration, Time Management, Agile/Scrum, Initiative, Documentation \\
  \textbf{Areas of Interest:} & Mobile App Development, Cross-platform Development, AI/ML, Research \\
\end{tabular}


\section{Selected Technical Projects}
\small
%%CAREEROS:PROJECTS%%


%-----------LEADERSHIP, INVOLVEMENT, AND CERTIFICATIONS-----------------
\section{Leadership, Involvement, and Certifications}
\small
\begin{itemize}[leftmargin=*, itemsep=1pt, topsep=1pt]

  \item \textbf{IEEE Student Branch, VIT Pune} \emph{(2021 -- 2024)}:
  Progressed from \textbf{Executive Committee Member} to \textbf{Project Head}. Contributed as part of the App Team, developing and maintaining multiple projects while mentoring peers in app development. Conducted internal coding sessions and organized hands-on workshops to strengthen technical knowledge within the branch.

  \item \textbf{App Team Leadership:}
  As Project Head, guided and mentored 35 executive committee members in Flutter, app development, and open-source practices. Structured a workshop-based learning model with weekly tasks, identified top performers, and formed a core app team to build the official IEEE SB mobile application. Supervised two additional industry projects under the team’s portfolio.

  \item \textbf{DevOps Workshop Instructor:}
  Served as a technical instructor for a large-scale DevOps workshop, delivering hands-on sessions to over 60 participants covering Docker, CI/CD, Kubernetes, and cloud deployment fundamentals.

  \item \textbf{Hacktoberfest Contributor \& Mentor} \emph{(2022 -- 2023)}:
  Achieved \textbf{Level 4 Contributor} for two consecutive years by making consistent open-source contributions. Later introduced and mentored IEEE members in contributing to open source, helping many achieve their first successful PRs.

  \item \textbf{Student Council, VIT Pune} \emph{(2023 -- 2024)}:
  Served as \textbf{Finance Coordinator}, managing event budgets, vendor coordination, and financial documentation for major campus events. Collaborated closely with faculty and management to ensure smooth execution of college-wide initiatives.

  \item \textbf{Certifications:}
  Earned the \textbf{IBM Full Stack Software Development Specialization (Coursera)} — a comprehensive 15-course program covering front-end, back-end, DevOps, and cloud deployment workflows. Additional certifications in Machine Learning, Deep Learning, and Data Analysis using Python.

  \item \textbf{Anthropic Certifications:}
  Completed 5 Anthropic certifications — \textbf{Introduction to Model Context Protocol}, \textbf{Claude with the Anthropic API}, \textbf{Introduction to Agent Skills}, \textbf{Claude Code in Action}, and \textbf{AI Fluency: Framework \& Foundations}.

\end{itemize}

\section{Education}
\begin{tabularx}{\linewidth}{@{}l X r@{}}
2021 -- 2025 & \textbf{B.Tech in Computer Engineering}, Vishwakarma Institute of Technology (VIT), Pune & \textit{CGPA: 8.64/10.0} \\
& \textbf{Relevant Coursework:} Data Structures and Algorithms, Operating Systems, Database Management Systems, Artificial Intelligence, Machine Learning, Computer Networks & \\
\\
2020 -- 2021 & \textbf{Higher Secondary Education (HSC)}, Deogiri College, Aurangabad & \textit{Percentage: 93.33\%} \\
2018 -- 2019 & \textbf{Class 10\textsuperscript{th}}, Stepping Stones High School, Aurangabad (CBSE) & \textit{Percentage: 87.60\%} \\
\end{tabularx}

\end{document}
